% https://www.latextemplates.com/template/the-legrand-orange-book

%%%%%%%%%%%%%%%%%%%%%%%%%%%%%%%%%%%%%%%%%
% The Legrand Orange Book
% LaTeX Template
% Version 2.4 (26/09/2018)
%
% This template was downloaded from:
% http://www.LaTeXTemplates.com
%
% Original author:
% Mathias Legrand (legrand.mathias@gmail.com) with modifications by:
% Vel (vel@latextemplates.com)
%
% License:
% CC BY-NC-SA 3.0 (http://creativecommons.org/licenses/by-nc-sa/3.0/)
%
% Compiling this template:
% This template uses biber for its bibliography and makeindex for its index.
% When you first open the template, compile it from the command line with the 
% commands below to make sure your LaTeX distribution is configured correctly:
%
% 1) pdflatex main
% 2) makeindex main.idx -s StyleInd.ist
% 3) biber main
% 4) pdflatex main x 2
%
% After this, when you wish to update the bibliography/index use the appropriate
% command above and make sure to compile with pdflatex several times 
% afterwards to propagate your changes to the document.
%
% This template also uses a number of packages which may need to be
% updated to the newest versions for the template to compile. It is strongly
% recommended you update your LaTeX distribution if you have any
% compilation errors.
%
% Important note:
% Chapter heading images should have a 2:1 width:height ratio,
% e.g. 920px width and 460px height.
%
%%%%%%%%%%%%%%%%%%%%%%%%%%%%%%%%%%%%%%%%%

%----------------------------------------------------------------------------------------
%	PACKAGES AND OTHER DOCUMENT CONFIGURATIONS
%----------------------------------------------------------------------------------------

\documentclass[11pt,fleqn]{book} % Default font size and left-justified equations

%%%%%%%%%%%%%%%%%%%%%%%%%%%%%%%%%%%%%%%%%
% The Legrand Orange Book
% Structural Definitions File
% Version 2.1 (26/09/2018)
%
% Original author:
% Mathias Legrand (legrand.mathias@gmail.com) with modifications by:
% Vel (vel@latextemplates.com)
% 
% This file was downloaded from:
% http://www.LaTeXTemplates.com
%
% License:
% CC BY-NC-SA 3.0 (http://creativecommons.org/licenses/by-nc-sa/3.0/)
%
%%%%%%%%%%%%%%%%%%%%%%%%%%%%%%%%%%%%%%%%%

%----------------------------------------------------------------------------------------
%	VARIOUS REQUIRED PACKAGES AND CONFIGURATIONS
%----------------------------------------------------------------------------------------

\usepackage{graphicx} % Required for including pictures
\graphicspath{{Pictures/}} % Specifies the directory where pictures are stored

\usepackage{lipsum} % Inserts dummy text

\usepackage{tikz} % Required for drawing custom shapes

\usepackage[portuguese]{babel} % English language/hyphenation

\usepackage{enumitem} % Customize lists
\setlist{nolistsep} % Reduce spacing between bullet points and numbered lists

\usepackage{booktabs} % Required for nicer horizontal rules in tables

\usepackage{xcolor} % Required for specifying colors by name
\definecolor{ocre}{RGB}{243,102,25} % Define the orange color used for highlighting throughout the book

%----------------------------------------------------------------------------------------
%	MARGINS
%----------------------------------------------------------------------------------------

\usepackage{geometry} % Required for adjusting page dimensions and margins

\geometry{
	paper=a4paper, % Paper size, change to letterpaper for US letter size
	top=3cm, % Top margin
	bottom=3cm, % Bottom margin
	left=3cm, % Left margin
	right=3cm, % Right margin
	headheight=14pt, % Header height
	footskip=1.4cm, % Space from the bottom margin to the baseline of the footer
	headsep=10pt, % Space from the top margin to the baseline of the header
	%showframe, % Uncomment to show how the type block is set on the page
}

%----------------------------------------------------------------------------------------
%	FONTS
%----------------------------------------------------------------------------------------

\usepackage{avant} % Use the Avantgarde font for headings
%\usepackage{times} % Use the Times font for headings
\usepackage{mathptmx} % Use the Adobe Times Roman as the default text font together with math symbols from the Sym­bol, Chancery and Com­puter Modern fonts

\usepackage{microtype} % Slightly tweak font spacing for aesthetics
\usepackage[utf8]{inputenc} % Required for including letters with accents
\usepackage[T1]{fontenc} % Use 8-bit encoding that has 256 glyphs

%----------------------------------------------------------------------------------------
%	BIBLIOGRAPHY AND INDEX
%----------------------------------------------------------------------------------------

\usepackage[style=numeric,citestyle=numeric,sorting=nyt,sortcites=true,autopunct=true,babel=hyphen,hyperref=true,abbreviate=false,backref=true,backend=biber]{biblatex}
\addbibresource{bibliography.bib} % BibTeX bibliography file
\defbibheading{bibempty}{}

\usepackage{calc} % For simpler calculation - used for spacing the index letter headings correctly
\usepackage{makeidx} % Required to make an index
\makeindex % Tells LaTeX to create the files required for indexing

%----------------------------------------------------------------------------------------
%	MAIN TABLE OF CONTENTS
%----------------------------------------------------------------------------------------

\usepackage{titletoc} % Required for manipulating the table of contents

\contentsmargin{0cm} % Removes the default margin

% Part text styling (this is mostly taken care of in the PART HEADINGS section of this file)
\titlecontents{part}
	[0cm] % Left indentation
	{\addvspace{20pt}\bfseries} % Spacing and font options for parts
	{}
	{}
	{}

% Chapter text styling
\titlecontents{chapter}
	[1.25cm] % Left indentation
	{\addvspace{12pt}\large\sffamily\bfseries} % Spacing and font options for chapters
	{\color{ocre!60}\contentslabel[\Large\thecontentslabel]{1.25cm}\color{ocre}} % Formatting of numbered sections of this type
	{\color{ocre}} % Formatting of numberless sections of this type
	{\color{ocre!60}\normalsize\;\titlerule*[.5pc]{.}\;\thecontentspage} % Formatting of the filler to the right of the heading and the page number

% Section text styling
\titlecontents{section}
	[1.25cm] % Left indentation
	{\addvspace{3pt}\sffamily\bfseries} % Spacing and font options for sections
	{\contentslabel[\thecontentslabel]{1.25cm}} % Formatting of numbered sections of this type
	{} % Formatting of numberless sections of this type
	{\hfill\color{black}\thecontentspage} % Formatting of the filler to the right of the heading and the page number

% Subsection text styling
\titlecontents{subsection}
	[1.25cm] % Left indentation
	{\addvspace{1pt}\sffamily\small} % Spacing and font options for subsections
	{\contentslabel[\thecontentslabel]{1.25cm}} % Formatting of numbered sections of this type
	{} % Formatting of numberless sections of this type
	{\ \titlerule*[.5pc]{.}\;\thecontentspage} % Formatting of the filler to the right of the heading and the page number

% Figure text styling
\titlecontents{figure}
	[1.25cm] % Left indentation
	{\addvspace{1pt}\sffamily\small} % Spacing and font options for figures
	{\thecontentslabel\hspace*{1em}} % Formatting of numbered sections of this type
	{} % Formatting of numberless sections of this type
	{\ \titlerule*[.5pc]{.}\;\thecontentspage} % Formatting of the filler to the right of the heading and the page number

% Table text styling
\titlecontents{table}
	[1.25cm] % Left indentation
	{\addvspace{1pt}\sffamily\small} % Spacing and font options for tables
	{\thecontentslabel\hspace*{1em}} % Formatting of numbered sections of this type
	{} % Formatting of numberless sections of this type
	{\ \titlerule*[.5pc]{.}\;\thecontentspage} % Formatting of the filler to the right of the heading and the page number

%----------------------------------------------------------------------------------------
%	MINI TABLE OF CONTENTS IN PART HEADS
%----------------------------------------------------------------------------------------

% Chapter text styling
\titlecontents{lchapter}
	[0em] % Left indentation
	{\addvspace{15pt}\large\sffamily\bfseries} % Spacing and font options for chapters
	{\color{ocre}\contentslabel[\Large\thecontentslabel]{1.25cm}\color{ocre}} % Chapter number
	{}  
	{\color{ocre}\normalsize\sffamily\bfseries\;\titlerule*[.5pc]{.}\;\thecontentspage} % Page number

% Section text styling
\titlecontents{lsection}
	[0em] % Left indentation
	{\sffamily\small} % Spacing and font options for sections
	{\contentslabel[\thecontentslabel]{1.25cm}} % Section number
	{}
	{}

% Subsection text styling (note these aren't shown by default, display them by searchings this file for tocdepth and reading the commented text)
\setcounter{tocdepth}{0}
\titlecontents{lsubsection}
	[.5em] % Left indentation
	{\sffamily\footnotesize} % Spacing and font options for subsections
	{\contentslabel[\thecontentslabel]{1.25cm}}
	{}
	{}

%----------------------------------------------------------------------------------------
%	HEADERS AND FOOTERS
%----------------------------------------------------------------------------------------

\usepackage{fancyhdr} % Required for header and footer configuration

\pagestyle{fancy} % Enable the custom headers and footers

\renewcommand{\chaptermark}[1]{\markboth{\sffamily\normalsize\bfseries\chaptername\ \thechapter.\ #1}{}} % Styling for the current chapter in the header
\renewcommand{\sectionmark}[1]{\markright{\sffamily\normalsize\thesection\hspace{5pt}#1}{}} % Styling for the current section in the header

\fancyhf{} % Clear default headers and footers
\fancyhead[LE,RO]{\sffamily\normalsize\thepage} % Styling for the page number in the header
\fancyhead[LO]{\rightmark} % Print the nearest section name on the left side of odd pages
\fancyhead[RE]{\leftmark} % Print the current chapter name on the right side of even pages
%\fancyfoot[C]{\thepage} % Uncomment to include a footer

\renewcommand{\headrulewidth}{0.5pt} % Thickness of the rule under the header

\fancypagestyle{plain}{% Style for when a plain pagestyle is specified
	\fancyhead{}\renewcommand{\headrulewidth}{0pt}%
}

% Removes the header from odd empty pages at the end of chapters
\makeatletter
\renewcommand{\cleardoublepage}{
\clearpage\ifodd\c@page\else
\hbox{}
\vspace*{\fill}
\thispagestyle{empty}
\newpage
\fi}

%----------------------------------------------------------------------------------------
%	THEOREM STYLES
%----------------------------------------------------------------------------------------

\usepackage{amsmath,amsfonts,amssymb,amsthm} % For math equations, theorems, symbols, etc

\newcommand{\intoo}[2]{\mathopen{]}#1\,;#2\mathclose{[}}
\newcommand{\ud}{\mathop{\mathrm{{}d}}\mathopen{}}
\newcommand{\intff}[2]{\mathopen{[}#1\,;#2\mathclose{]}}
\renewcommand{\qedsymbol}{$\blacksquare$}
\newtheorem{notation}{Notation}[chapter]

% Boxed/framed environments
\newtheoremstyle{ocrenumbox}% Theorem style name
{0pt}% Space above
{0pt}% Space below
{\normalfont}% Body font
{}% Indent amount
{\small\bf\sffamily\color{ocre}}% Theorem head font
{\;}% Punctuation after theorem head
{0.25em}% Space after theorem head
{\small\sffamily\color{ocre}\thmname{#1}\nobreakspace\thmnumber{\@ifnotempty{#1}{}\@upn{#2}}% Theorem text (e.g. Theorem 2.1)
\thmnote{\nobreakspace\the\thm@notefont\sffamily\bfseries\color{black}---\nobreakspace#3.}} % Optional theorem note

\newtheoremstyle{blacknumex}% Theorem style name
{5pt}% Space above
{5pt}% Space below
{\normalfont}% Body font
{} % Indent amount
{\small\bf\sffamily}% Theorem head font
{\;}% Punctuation after theorem head
{0.25em}% Space after theorem head
{\small\sffamily{\tiny\ensuremath{\blacksquare}}\nobreakspace\thmname{#1}\nobreakspace\thmnumber{\@ifnotempty{#1}{}\@upn{#2}}% Theorem text (e.g. Theorem 2.1)
\thmnote{\nobreakspace\the\thm@notefont\sffamily\bfseries---\nobreakspace#3.}}% Optional theorem note

\newtheoremstyle{blacknumbox} % Theorem style name
{0pt}% Space above
{0pt}% Space below
{\normalfont}% Body font
{}% Indent amount
{\small\bf\sffamily}% Theorem head font
{\;}% Punctuation after theorem head
{0.25em}% Space after theorem head
{\small\sffamily\thmname{#1}\nobreakspace\thmnumber{\@ifnotempty{#1}{}\@upn{#2}}% Theorem text (e.g. Theorem 2.1)
\thmnote{\nobreakspace\the\thm@notefont\sffamily\bfseries---\nobreakspace#3.}}% Optional theorem note

% Non-boxed/non-framed environments
\newtheoremstyle{ocrenum}% Theorem style name
{5pt}% Space above
{5pt}% Space below
{\normalfont}% Body font
{}% Indent amount
{\small\bf\sffamily\color{ocre}}% Theorem head font
{\;}% Punctuation after theorem head
{0.25em}% Space after theorem head
{\small\sffamily\color{ocre}\thmname{#1}\nobreakspace\thmnumber{\@ifnotempty{#1}{}\@upn{#2}}% Theorem text (e.g. Theorem 2.1)
\thmnote{\nobreakspace\the\thm@notefont\sffamily\bfseries\color{black}---\nobreakspace#3.}} % Optional theorem note
\makeatother

% Defines the theorem text style for each type of theorem to one of the three styles above
\newcounter{dummy} 
\numberwithin{dummy}{section}
\theoremstyle{ocrenumbox}
\newtheorem{theoremeT}[dummy]{Theorem}
\newtheorem{problem}{Problem}[chapter]
\newtheorem{exerciseT}{Exercise}[chapter]
\theoremstyle{blacknumex}
\newtheorem{exampleT}{Example}[chapter]
\theoremstyle{blacknumbox}
\newtheorem{vocabulary}{Vocabulary}[chapter]
\newtheorem{definitionT}{Definition}[section]
\newtheorem{corollaryT}[dummy]{Corollary}
\theoremstyle{ocrenum}
\newtheorem{proposition}[dummy]{Proposition}

%----------------------------------------------------------------------------------------
%	DEFINITION OF COLORED BOXES
%----------------------------------------------------------------------------------------

\RequirePackage[framemethod=default]{mdframed} % Required for creating the theorem, definition, exercise and corollary boxes

% Theorem box
\newmdenv[skipabove=7pt,
skipbelow=7pt,
backgroundcolor=black!5,
linecolor=ocre,
innerleftmargin=5pt,
innerrightmargin=5pt,
innertopmargin=5pt,
leftmargin=0cm,
rightmargin=0cm,
innerbottommargin=5pt]{tBox}

% Exercise box	  
\newmdenv[skipabove=7pt,
skipbelow=7pt,
rightline=false,
leftline=true,
topline=false,
bottomline=false,
backgroundcolor=ocre!10,
linecolor=ocre,
innerleftmargin=5pt,
innerrightmargin=5pt,
innertopmargin=5pt,
innerbottommargin=5pt,
leftmargin=0cm,
rightmargin=0cm,
linewidth=4pt]{eBox}	

% Definition box
\newmdenv[skipabove=7pt,
skipbelow=7pt,
rightline=false,
leftline=true,
topline=false,
bottomline=false,
linecolor=ocre,
innerleftmargin=5pt,
innerrightmargin=5pt,
innertopmargin=0pt,
leftmargin=0cm,
rightmargin=0cm,
linewidth=4pt,
innerbottommargin=0pt]{dBox}	

% Corollary box
\newmdenv[skipabove=7pt,
skipbelow=7pt,
rightline=false,
leftline=true,
topline=false,
bottomline=false,
linecolor=gray,
backgroundcolor=black!5,
innerleftmargin=5pt,
innerrightmargin=5pt,
innertopmargin=5pt,
leftmargin=0cm,
rightmargin=0cm,
linewidth=4pt,
innerbottommargin=5pt]{cBox}

% Creates an environment for each type of theorem and assigns it a theorem text style from the "Theorem Styles" section above and a colored box from above
\newenvironment{theorem}{\begin{tBox}\begin{theoremeT}}{\end{theoremeT}\end{tBox}}
\newenvironment{exercise}{\begin{eBox}\begin{exerciseT}}{\hfill{\color{ocre}\tiny\ensuremath{\blacksquare}}\end{exerciseT}\end{eBox}}				  
\newenvironment{definition}{\begin{dBox}\begin{definitionT}}{\end{definitionT}\end{dBox}}	
\newenvironment{example}{\begin{exampleT}}{\hfill{\tiny\ensuremath{\blacksquare}}\end{exampleT}}		
\newenvironment{corollary}{\begin{cBox}\begin{corollaryT}}{\end{corollaryT}\end{cBox}}	

%----------------------------------------------------------------------------------------
%	REMARK ENVIRONMENT
%----------------------------------------------------------------------------------------

\newenvironment{remark}{\par\vspace{10pt}\small % Vertical white space above the remark and smaller font size
\begin{list}{}{
\leftmargin=35pt % Indentation on the left
\rightmargin=25pt}\item\ignorespaces % Indentation on the right
\makebox[-2.5pt]{\begin{tikzpicture}[overlay]
\node[draw=ocre!60,line width=1pt,circle,fill=ocre!25,font=\sffamily\bfseries,inner sep=2pt,outer sep=0pt] at (-15pt,0pt){\textcolor{ocre}{R}};\end{tikzpicture}} % Orange R in a circle
\advance\baselineskip -1pt}{\end{list}\vskip5pt} % Tighter line spacing and white space after remark

%----------------------------------------------------------------------------------------
%	SECTION NUMBERING IN THE MARGIN
%----------------------------------------------------------------------------------------

\makeatletter
\renewcommand{\@seccntformat}[1]{\llap{\textcolor{ocre}{\csname the#1\endcsname}\hspace{1em}}}                    
\renewcommand{\section}{\@startsection{section}{1}{\z@}
{-4ex \@plus -1ex \@minus -.4ex}
{1ex \@plus.2ex }
{\normalfont\large\sffamily\bfseries}}
\renewcommand{\subsection}{\@startsection {subsection}{2}{\z@}
{-3ex \@plus -0.1ex \@minus -.4ex}
{0.5ex \@plus.2ex }
{\normalfont\sffamily\bfseries}}
\renewcommand{\subsubsection}{\@startsection {subsubsection}{3}{\z@}
{-2ex \@plus -0.1ex \@minus -.2ex}
{.2ex \@plus.2ex }
{\normalfont\small\sffamily\bfseries}}                        
\renewcommand\paragraph{\@startsection{paragraph}{4}{\z@}
{-2ex \@plus-.2ex \@minus .2ex}
{.1ex}
{\normalfont\small\sffamily\bfseries}}

%----------------------------------------------------------------------------------------
%	PART HEADINGS
%----------------------------------------------------------------------------------------

% Numbered part in the table of contents
\newcommand{\@mypartnumtocformat}[2]{%
	\setlength\fboxsep{0pt}%
	\noindent\colorbox{ocre!20}{\strut\parbox[c][.7cm]{\ecart}{\color{ocre!70}\Large\sffamily\bfseries\centering#1}}\hskip\esp\colorbox{ocre!40}{\strut\parbox[c][.7cm]{\linewidth-\ecart-\esp}{\Large\sffamily\centering#2}}%
}

% Unnumbered part in the table of contents
\newcommand{\@myparttocformat}[1]{%
	\setlength\fboxsep{0pt}%
	\noindent\colorbox{ocre!40}{\strut\parbox[c][.7cm]{\linewidth}{\Large\sffamily\centering#1}}%
}

\newlength\esp
\setlength\esp{4pt}
\newlength\ecart
\setlength\ecart{1.2cm-\esp}
\newcommand{\thepartimage}{}%
\newcommand{\partimage}[1]{\renewcommand{\thepartimage}{#1}}%
\def\@part[#1]#2{%
\ifnum \c@secnumdepth >-2\relax%
\refstepcounter{part}%
\addcontentsline{toc}{part}{\texorpdfstring{\protect\@mypartnumtocformat{\thepart}{#1}}{\partname~\thepart\ ---\ #1}}
\else%
\addcontentsline{toc}{part}{\texorpdfstring{\protect\@myparttocformat{#1}}{#1}}%
\fi%
\startcontents%
\markboth{}{}%
{\thispagestyle{empty}%
\begin{tikzpicture}[remember picture,overlay]%
\node at (current page.north west){\begin{tikzpicture}[remember picture,overlay]%	
\fill[ocre!20](0cm,0cm) rectangle (\paperwidth,-\paperheight);
\node[anchor=north] at (4cm,-3.25cm){\color{ocre!40}\fontsize{220}{100}\sffamily\bfseries\thepart}; 
\node[anchor=south east] at (\paperwidth-1cm,-\paperheight+1cm){\parbox[t][][t]{8.5cm}{
\printcontents{l}{0}{\setcounter{tocdepth}{1}}% The depth to which the Part mini table of contents displays headings; 0 for chapters only, 1 for chapters and sections and 2 for chapters, sections and subsections
}};
\node[anchor=north east] at (\paperwidth-1.5cm,-3.25cm){\parbox[t][][t]{15cm}{\strut\raggedleft\color{white}\fontsize{30}{30}\sffamily\bfseries#2}};
\end{tikzpicture}};
\end{tikzpicture}}%
\@endpart}
\def\@spart#1{%
\startcontents%
\phantomsection
{\thispagestyle{empty}%
\begin{tikzpicture}[remember picture,overlay]%
\node at (current page.north west){\begin{tikzpicture}[remember picture,overlay]%	
\fill[ocre!20](0cm,0cm) rectangle (\paperwidth,-\paperheight);
\node[anchor=north east] at (\paperwidth-1.5cm,-3.25cm){\parbox[t][][t]{15cm}{\strut\raggedleft\color{white}\fontsize{30}{30}\sffamily\bfseries#1}};
\end{tikzpicture}};
\end{tikzpicture}}
\addcontentsline{toc}{part}{\texorpdfstring{%
\setlength\fboxsep{0pt}%
\noindent\protect\colorbox{ocre!40}{\strut\protect\parbox[c][.7cm]{\linewidth}{\Large\sffamily\protect\centering #1\quad\mbox{}}}}{#1}}%
\@endpart}
\def\@endpart{\vfil\newpage
\if@twoside
\if@openright
\null
\thispagestyle{empty}%
\newpage
\fi
\fi
\if@tempswa
\twocolumn
\fi}

%----------------------------------------------------------------------------------------
%	CHAPTER HEADINGS
%----------------------------------------------------------------------------------------

% A switch to conditionally include a picture, implemented by Christian Hupfer
\newif\ifusechapterimage
\usechapterimagetrue
\newcommand{\thechapterimage}{}%
\newcommand{\chapterimage}[1]{\ifusechapterimage\renewcommand{\thechapterimage}{#1}\fi}%
\newcommand{\autodot}{.}
\def\@makechapterhead#1{%
{\parindent \z@ \raggedright \normalfont
\ifnum \c@secnumdepth >\m@ne
\if@mainmatter
\begin{tikzpicture}[remember picture,overlay]
\node at (current page.north west)
{\begin{tikzpicture}[remember picture,overlay]
\node[anchor=north west,inner sep=0pt] at (0,0) {\ifusechapterimage\includegraphics[width=\paperwidth]{\thechapterimage}\fi};
\draw[anchor=west] (\Gm@lmargin,-9cm) node [line width=2pt,rounded corners=15pt,draw=ocre,fill=white,fill opacity=0.5,inner sep=15pt]{\strut\makebox[22cm]{}};
\draw[anchor=west] (\Gm@lmargin+.3cm,-9cm) node {\huge\sffamily\bfseries\color{black}\thechapter\autodot~#1\strut};
\end{tikzpicture}};
\end{tikzpicture}
\else
\begin{tikzpicture}[remember picture,overlay]
\node at (current page.north west)
{\begin{tikzpicture}[remember picture,overlay]
\node[anchor=north west,inner sep=0pt] at (0,0) {\ifusechapterimage\includegraphics[width=\paperwidth]{\thechapterimage}\fi};
\draw[anchor=west] (\Gm@lmargin,-9cm) node [line width=2pt,rounded corners=15pt,draw=ocre,fill=white,fill opacity=0.5,inner sep=15pt]{\strut\makebox[22cm]{}};
\draw[anchor=west] (\Gm@lmargin+.3cm,-9cm) node {\huge\sffamily\bfseries\color{black}#1\strut};
\end{tikzpicture}};
\end{tikzpicture}
\fi\fi\par\vspace*{270\p@}}}

%-------------------------------------------

\def\@makeschapterhead#1{%
\begin{tikzpicture}[remember picture,overlay]
\node at (current page.north west)
{\begin{tikzpicture}[remember picture,overlay]
\node[anchor=north west,inner sep=0pt] at (0,0) {\ifusechapterimage\includegraphics[width=\paperwidth]{\thechapterimage}\fi};
\draw[anchor=west] (\Gm@lmargin,-9cm) node [line width=2pt,rounded corners=15pt,draw=ocre,fill=white,fill opacity=0.5,inner sep=15pt]{\strut\makebox[22cm]{}};
\draw[anchor=west] (\Gm@lmargin+.3cm,-9cm) node {\huge\sffamily\bfseries\color{black}#1\strut};
\end{tikzpicture}};
\end{tikzpicture}
\par\vspace*{270\p@}}
\makeatother

%----------------------------------------------------------------------------------------
%	LINKS
%----------------------------------------------------------------------------------------

\usepackage{hyperref}
\hypersetup{hidelinks,backref=true,pagebackref=true,hyperindex=true,colorlinks=false,breaklinks=true,urlcolor=ocre,bookmarks=true,bookmarksopen=false}

\usepackage{bookmark}
\bookmarksetup{
open,
numbered,
addtohook={%
\ifnum\bookmarkget{level}=0 % chapter
\bookmarksetup{bold}%
\fi
\ifnum\bookmarkget{level}=-1 % part
\bookmarksetup{color=ocre,bold}%
\fi
}
}

\hypersetup{
    colorlinks=true,
    linkcolor=blue,
    filecolor=magenta,      
    urlcolor=cyan,
    pdftitle={Overleaf Example},
    pdfpagemode=FullScreen,
    }


%\hypersetup{pdftitle={Title},pdfauthor={Author}} % Uncomment and fill out to include PDF metadata for the author and title of the book

%----------------------------------------------------------------------------------------

\begin{document}

%----------------------------------------------------------------------------------------
%	TITLE PAGE
%----------------------------------------------------------------------------------------

\begingroup
\thispagestyle{empty} % Suppress headers and footers on the title page
\begin{tikzpicture}[remember picture,overlay]
\node[inner sep=0pt] (background) at (current page.center) {\includegraphics[width=\paperwidth]{background.pdf}};
\draw (current page.center) node [fill=ocre!30!white,fill opacity=0.6,text opacity=1,inner sep=1cm]{\Huge\centering\bfseries\sffamily\parbox[c][][t]{\paperwidth}{\centering NEMO 2021\\[15pt] % Book title
{\Large Problemas e soluções}\\[20pt] % Subtitle
}}; % Author name
\end{tikzpicture}
\vfill
\endgroup

%----------------------------------------------------------------------------------------
%	COPYRIGHT PAGE
%----------------------------------------------------------------------------------------

\newpage
~\vfill
\thispagestyle{empty}

{\noindent\huge \href{https://t.me/joinchat/SPRsMWjgBF1SssVy}{Grupo no Telegram.}

\noindent\href{https://www.overleaf.com/8635267617vdpxqmfzjjry}{Link para este documento.}

\noindent\href{https://meet.google.com/gwn-qqyh-dvs}{Encontros aos sábados 18:00.}

\noindent\href{https://github.com/SofiaSL/NEMO}{Backup do \LaTeX\hphantom{ } no github.}\vspace{50px}


\large\noindent Para adicionar problemas ou soluções vá para o arquivo problemasN.tex ou solucoesN.tex onde N é a semana atual. A ideia é que você adicione problemas qualquer dia da semana, e o grupo escreva a solução na reunião. Mas se preferir pode escrever a solução antes do começo da reunião para poupar tempo. Justamente para isso que os problemas e soluções estão em capítulos diferentes. Se tiver recursos legais para compartilhar também pode adiciona-los ao apêndice, capitulo de sugestões de materiais.

}

%----------------------------------------------------------------------------------------
%	TABLE OF CONTENTS
%----------------------------------------------------------------------------------------

%\usechapterimagefalse % If you don't want to include a chapter image, use this to toggle images off - it can be enabled later with \usechapterimagetrue

\chapterimage{chapter_head_1.pdf} % Table of contents heading image

\pagestyle{empty} % Disable headers and footers for the following pages

\tableofcontents % Print the table of contents itself

\cleardoublepage % Forces the first chapter to start on an odd page so it's on the right side of the book

\pagestyle{fancy} % Enable headers and footers again

%----------------------------------------------------------------------------------------
%	PART
%----------------------------------------------------------------------------------------

\part{Problemas}

\chapterimage{chapter_head_2.pdf} % Chapter heading image

\chapter{Semana 1 (29 de maio)}
\section{OBM 2020 problema 1}

Defina como $n(R)$ a quantidade de maneiras de escrever $R=2x^2+3y^2+5z^2$ com $x,y,z\in\mathbb{Z}$. Qual é o valor de

$$\lim_{R\to\infty}\frac{n(1)+n(2)+\dots+n(R)}{R^{3/2}}?$$

\section{OBM 2019 - 1ª fase - problema 5}

Calcule

$$\sum_{n=1}^\infty\frac{2}{4n^2+8n+3}.$$

\section{O incrível problema dos reis caminhantes}

Em um tabuleiro estão alguns reis. A cada segundo, um rei qualquer se move para uma casa adjacente. Após uma quantidade finita de segundos, todos os reis percorreram cada casa do tabuleiro exatamente uma vez e voltaram para a sua casa inicial. Prove que houve um momento em que nenhum rei estava na sua casa inicial.

\chapter{Semana 2 (5 de junho)}
\section{Problema 1 SiM 2020}

Todo participante de um torneio joga com cada um dos outros jogadores exatamente uma vez. Após o torneio cada jogador faz uma lista com os nomes de todos os jogadores vencidos por ele e de todos os jogadores que foram vencidos pelos jogadores que ele venceu. Sabendo que nesse torneio não há empates, prove que existe um jogador cuja lista possui o nome de todos os outros jogadores.

\section{Problema 1.1.5 Additive Combinatorics, T. Tao and V. Vu}

Sejam $A$ e $B$ subconjuntos não vazios de um grupo aditivo finito $Z$. Mostre que existem $x,y\in Z$ tais que
$$1-\frac{|A\cup (B+x)|}{|Z|}\le \left(1-\frac{|A|}{|Z|}\right) \left(1-\frac{|B|}{|Z|}\right)\le 1-\frac{|A\cup (B+y)|}{|Z|}$$

\chapter{Semana 3 (12 de junho)}
\section{MIC OCW - Putnam Seminar}

Seja $p_n$ a probabilidade de que $c+d$ seja um quadrado perfeito com $c$ e $d$ escolhidos uniforme e independentemente no conjunto $\{1,\dots,n\}$. Mostre que $\lim_{n\to \infty} p_n\sqrt{n}$ existe e calcule seu valor.

\section{OBM 2001 (Láááá na pré-história)}

Seja $x_n$ uma sequência definida por
\begin{align}
    x_{n+1} = x^2_{n} - \frac{x_n}{2},\; n \ge 0
\end{align}
Para quais valores de $x_0$ a sequência converge? Para que valor?

\chapter{Semana 4 (19 de junho)}
\section{Putnam 2005}

Calcule

$$\int_0^1\frac{\ln(x+1)}{x^2+1}dx.$$

\section{Revista do Professor de Matemática 103 problema 454}

Na figura os dois círculos tem o mesmo raio e são tangentes aos lados do triângulo e às retas. O triângulo é equilátero de lado $2$. Qual é o raio dos círculos?

\begin{center}
\begin{tikzpicture}[scale=3]

\draw (-1,0) -- (1,0) -- (0, 1.73205) -- (-1,0);
\draw (0,1.09637) circle (0.317837);
\draw (0,0.31783) circle (0.317837);
\draw (-1,0) -- (0.42020,1.0042);
\draw (1,0) -- (-0.42020,1.0042);

\end{tikzpicture}
\end{center}

\section{Putnam and Beyond 1.1.3 (generalizado)}

Sejam $k$ inteiro positivo e $S = \{a_1,a_2,\;\dots,a_n\}$ um conjunto de inteiros positivos primos dois a dois tal que $1<a_i\le k,\; \forall a_i\in S$. Determine o menor $n$ para o qual $S$ contém ao menos um número primo.

\section{Revista do Professor de Matemática 103 problema 453/Dissertação lokona (generalizado)}

Seja $p(x) = x^n + a_{n-1} x^{n-1}+\;\dots+a_1x+a_0$ um polinômio de n raízes reais em que $a_{n-1} = -nk$, para $k$ algum real positivo. Sabendo que $p(k+1) > 1$, prove que $p(x)$ tem ao menos uma raiz maior que $k+1$. 

\section{Revista do Professor de Matemática 103 problema 452 (generalizado)}

Ache os inteiros $n$ tal que $\frac n2 +1$ são quadrados perfeitos.

\part{Soluções}

\chapter{Semana 1 (29 de maio)}
\section{OBM 2020 problema 1}

Esse problema é análogo ao fato que a função $f(n)$ que conta as maneiras de escrever $n=x^2+y^2$ com $x,y$ inteiros, tem limite

$$\frac{f(1)+f(2)+\dots+f(n)}{n}\to\pi$$

quando $n\to\infty$, isso por que o numerador é o número de pontos no plano com ambas as coordenadas inteiras contidos no circulo de raio $\sqrt{n}$ centrado na origem. Essa quantidade se aproxima da área do circulo e portanto o limite é $\pi$. Da mesma maneira, o limite da questão é dado pelo volume do elipsoide $2x^2+3y^2+5z^2=R$, que é

$$\frac{4}{3}\pi\sqrt{2}\sqrt{3}\sqrt{5}R^{3/2}.$$

\section{OBM 2019 problema 5}

Com a expansão em frações parciais:

$$\frac{2}{4n^2+8n+3}=\frac{1}{2n+1}-\frac{1}{2n+3}$$

Fica claro que a soma é telescópica e só o primeiro termo não cancela.

\section{O incrível problema dos reis caminhantes}



\chapter{Semana 2 (5 de junho)}
\section{Problema 1 SiM 2020}

Note que, se houver um jogador cuja lista possui o nome de todos os outros, tal lista deve ser a maior (no sentido de não haver outra com mais nomes) dentre a de todos os jogadores. Provaremos que qualquer lista com a maior quantidade de nomes tem o nome de todos os outros jogadores.\\

Seja $A$ a lista de Arnaldo, que possui a maior quantidade de nomes dentre as listas. Suponha, por absurdo, que exista um jogador, Bernardo, que não está na lista $A$; seja $B$ a lista de Bernardo. Assim, Bernardo venceu Arnaldo e todos os competidores que Arnaldo venceu. Em particular, $B$ tem mais nomes que $A$, \textbf{absurdo!} Portanto $A$ tem os nomes de todos os jogadores exceto Arnaldo.

\section{Problema 1.1.5 Additive Combinatorics, T. Tao and V. Vu}

Pelo princípio de inclusão-exclusão, é fato que:

$$1-\frac{|A\cup (B+x)|}{|Z|}\le \left(1-\frac{|A|}{|Z|}\right) \left(1-\frac{|B|}{|Z|}\right)\le 1-\frac{|A\cup (B+y)|}{|Z|}\; \Leftrightarrow$$

$$\Leftrightarrow 1 - \frac{|A|+|B+x|-|A \cap (B+x)|}{|Z|}\le 1 - \frac{|A|}{|Z|} - \frac{|B|}{|Z|} + \frac{|A||B|}{|Z|^2}\le 1 - \frac{|A|+|B+y|-|A \cap (B+y)|}{|Z|} \Leftrightarrow$$

$$\Leftrightarrow \; -\frac{|A|+|B|-|A \cap (B+x)|}{|Z|}\le - \frac{|A|}{|Z|} - \frac{|B|}{|Z|} + \frac{|A||B|}{|Z|^2}\le -\frac{|A|+|B|-|A \cap (B+y)|}{|Z|}\; \Leftrightarrow$$

$$\Leftrightarrow \; \frac{|A\cap (B+x)|}{|Z|}\le \frac{|A||B|}{|Z|^2}\le \frac{|A\cap (B+y)|}{|Z|} \; \Leftrightarrow$$

$$\Leftrightarrow \; |A\cap (B+x)||Z|\le |A||B|\le |A\cap (B+y)||Z| \hspace{1cm} (6.1)$$
\newline
Portanto, devemos provar que para quaisquer subconjuntos $A$ e $B$ de $Z$, existem $x,y\in Z$ que satisfazem (6.1).
\newline \newline
\textbf{Solução 1 (by Arjuna):}
\newline
Seja $B(z) := B+z$. Note que existem |Z| conjuntos (não necessariamente distintos) da forma $B+z,z\in Z$. Nesse sentido, note que cada $w\in Z$ pertence a $|B|$ desses conjuntos, pois para cada $b\in B$, temos que $w\in B(w-b)$. Assim, os elementos de $A$ aparecem $|A||B|$ vezes no total nos conjuntos $B(z),z\in Z$. Pelo princípio das casas dos pombos, pelo menos um desses conjuntos contém $\frac{|A||B|}{|Z|}$ elementos de $A$, ou seja:
$$\frac{|A||B|}{|Z|}\le |A\cap (B+y)|\hspace{0.3cm} \text{para algum $y\in Z$}$$
$$\Leftrightarrow \hspace{0.3cm} |A||B|\le |A\cap (B+y)||Z|\hspace{0.3cm} \text{para algum $y\in Z$}$$
Ademais, ao menos um desses conjuntos contém no máximo $\frac{|A||B|}{|Z|}$ elementos de $A$ (prova-se-lo trivialmente por absurdo), em outras palavras:
$$|A\cap (B+x)|\le \frac{|A||B|}{|Z|}\hspace{0.3cm} \text{para algum $x\in Z$}$$
$$\Leftrightarrow \hspace{0.3cm}|A\cap (B+x)||Z|\le |A||B|\hspace{0.3cm} \text{para algum $x\in Z$}$$
Como queríamos demonstrar.
\newline \newline
\textbf{Solução 2 (by Libardi):}
\newline
Seja $\Omega =\mathcal{P}(Z) = \{A|A\subseteq Z\}$. Nesse sentido, defina uma ação direita de $(Z,+)$ sobre $\Omega$ como $f:Z\times \Omega \rightarrow \Omega /f(x, A) = A+x$. Verificar que $f$ obedece aos axiomas de ação direita fica a cargo do leitor, mas vale relembrar a relação $|G_B||\mathcal{O}_B| = |Z|$, no qual $G_B$ é o grupo estabilizador de $B$ e $\mathcal{O}_B$ é a órbita de $B$. Nesse sentido, note que cada elemento $z\in Z$ pertence a $\frac{|B|}{|G_B|}$ conjuntos distintos de $\mathcal{O}_B$, pois para cada $w\in Z/z\in(B+w)$, temos que $(B+w)=(B+g)+w$ para todo $g\in G_B$. Sendo assim, os elementos de $A$ aparecem $\frac{|A||B|}{|G_B|}$ vezes nos conjuntos de $\mathcal{O}_B$. Pelo princípio da casa dos pombos, ao menos um dos conjuntos que pertence à $\mathcal{O}_B$ tem $\frac{|A||B|}{|G_B||\mathcal{O}_B|}$ elementos de $A$, ou seja:
$$\frac{|A||B|}{|G_B||\mathcal{O}_B|}\le |A\cap (B+y)|\hspace{0.3cm} \text{para algum $y\in Z$}$$
$$\Leftrightarrow \hspace{0.3cm} |A||B|\le |A\cap (B+y)||Z|\hspace{0.3cm} \text{para algum $y\in Z$}$$
Ademais, ao menos um desses conjuntos contém no máximo $\frac{|A||B|}{|G_B||\mathcal{O}_B|}$ elementos de $A$ (prova-se-lo trivialmente por absurdo), em outras palavras:
$$|A\cap (B+x)|\le \frac{|A||B|}{|G_B||\mathcal{O}_B|}\hspace{0.3cm} \text{para algum $x\in Z$}$$
$$\Leftrightarrow \hspace{0.3cm}|A\cap (B+x)||Z|\le |A||B|\hspace{0.3cm} \text{para algum $x\in Z$}$$
Como queríamos demonstrar.

\chapter{Semana 3 (12 de junho)}
\chapter{}

Note que se $(x_n)$ converge para $L\in \mathbb{R}$, então $L=L^2-\frac{L}{2}\Rightarrow L\in\{0,\frac{3}{2}\}$. Além disso, seja $f(x)=x^2-\frac{x}{2}$ para todo $x\in \mathbb{R}$; temos $x_{n+1}=f(x_n)$ para cada $n\ge 0$.\\

1) Se $x_0\in \{-1,\frac{3}{2}\},\ x_n=\frac{3}{2}\ \forall n\ge 1\Rightarrow \lim x_n=\frac{3}{2}$.\\

2) Se $x_0\in (-\infty,-1)\cup (\frac{3}{2},+\infty)$, então $x_1>\frac{3}{2}$, e como $g(x):=x-f(x)=\frac{3}{2}x-x^2<0$ para todo $x>\frac{3}{2}$,temos $x_n\ge x_1>\frac{3}{2}\ \forall n\ge 1$, donde a sequência $(x_n)$ diverge.\\

3) Se $x_0\in \left[0,\frac{3}{4}\right]$, note que $X:=f\left(\left[0,\frac{3}{4}\right]\right)=\left[-\frac{1}{16},\frac{3}{16}\right]$ é fechado (e portanto espaço métrico completo com a distância usual na reta) e tal que $f(X)\subset X$. Como $-\frac{5}{8}\le f'(x)\le -\frac{1}{8}\ \forall x\in X$, aplicando o TEorema do Valor Médio temos
$$|f(a)-f(b)|\le \frac{5}{8}|a-b|\ \forall a,b\in X$$
Pelo Teorema do Ponto Fixo de Banach, $(x_n)=(f^{n-1}(x_1))$ converge para o único ponto fixo de $f$ em $X$, i.e., $\lim x_n=0$.\\

4) Se $x_0\in \left(\frac{3}{4},\frac{3}{2}\right)$, observe que  $x_1\in \left[0,\frac{3}{2}\right]$ e que $g_{|\left(\frac{3}{4},\frac{3}{2}\right)}$ tem as seguintes propriedades:\\
a) $g_{|\left(\frac{3}{4},\frac{3}{2}\right)}>0$;\\
b) $g_{|\left(\frac{3}{4},\frac{3}{2}\right)}$ é estritamente decrescente
Suponha que $x_n\in \left(\frac{3}{4},\frac{3}{2}\right)$ para todo $n\ge 0$; então $\frac{3}{4}<x_n\le x_0<\frac{3}{2}$ e $g(x_n)\ge g(x_0)>0$ para todo $n\ge 0$. Portanto:
$$x_0-x_n=\sum_{i=0}^{n-1} x_i-x_{i+1}=\sum_{i=0}^{n-1} x_i-f(x_i)=\sum_{i=0}^{n-1} g(x_i)\ge ng(x_0)\overset{n\to \infty}{\rightarrow} \infty$$
$\textbf{ABSURDO!}$. Logo $x_k\in \left[0,\frac{3}{4}\right]$ para algum $k\ge 0$, donde $(x_n)=(f^{n-k})(x_k)$ converge para $0$ por $3$.

5) Se $x_0\in (-1,0)\Rightarrow x_1\in \left(0,\frac{3}{2}\right)$, e por $3)$ e $4)$, $\lim x_n=0$.

\chapter{Semana 4 (19 de junho)}
\section{Putnam 2005}
\textbf{Solução 1(by Sofia):}
\newline
Note que
$$\int_0^1\frac{\ln(x+1)}{x^2+1}dx=\int_0^1\sum_{n=1}^\infty \frac{(-1)^{n+1}}{n}\frac{x^n}{x^2+1}dx$$

$$=\sum_{n=1}^\infty\frac{(-1)^{n+1}}{n}\int_0^1 \frac{x^n}{x^2+1}dx$$

Defina

$$J_n=\int_0^1 \frac{x^n}{x^2+1}dx$$

$$\int_0^1\frac{\ln(x+1)}{x^2+1}dx=\sum_{n=1}^\infty\frac{(-1)^{n+1}}{n}J_n$$

então

$$J_n+J_{n+2}=\int_0^1 x^n dx=\frac{1}{n+1}.$$

$$J_0=\frac\pi 4$$

$$J_1=\frac{\log(2)}{2}$$

$$J_{2n}=(-1)^n\left(\frac\pi 4-\sum_{h=1}^n\frac{(-1)^{h+1}}{2h-1}\right)$$

$$J_{2n+1}=(-1)^n\left(\frac{\log(2)}{2}-\sum_{h=1}^n\frac{(-1)^{h+1}}{2h}\right)$$

\newline
\textbf{Solução 2 (by Libardi):}
\newline
Defina a função $I:\mathbb{R}\rightarrow\mathbb{R}$ de tal forma que:

$$I(t) = \int_0^t\frac{\ln(tx+1)}{x^2+1}dx$$

Trivialmente, temos $I(0) = 0$ e:

$$\Rightarrow \; \frac{dI}{dt} = \int_0^t\frac{x}{(tx+1)(x^2+1)}dx+\frac{\ln(t^2+1)}{t^2+1}$$

Fazendo frações parciais, determina-se $\frac{dI}{dt}$:

$$\frac{dI}{dt} = \frac{1}{t^2+1}\int_0^t\left(\frac{x+t}{x^2+1}-\frac{t}{tx+1}\right)dx+\frac{\ln(t^2+1)}{t^2+1}$$

$$\Leftrightarrow \frac{dI}{dt} = \frac{\ln(x^2+1)+2t\arctan(x)-2\ln(tx+1)}{2(t^2+1)}\Big|_{x=0}^{x=t} + \frac{\ln(t^2+1)}{t^2+1}$$

$$\Leftrightarrow \frac{dI}{dt} = \frac{t\arctan(x)}{t^2+1} + \frac{\ln(t^2+1)}{2(t^2+1)}$$

Para obter a integral desejada, integra-se $I'(t)$:

$$\int \frac{dI}{dt}dt = \int\left(\frac{t\arctan(x)}{t^2+1} + \frac{\ln(t^2+1)}{2(t^2+1)}\right)dt = \frac{1}{2}\arctan(t)\ln(1+t^2) + C$$

Portanto, como temos a condição de contorno $I(0) = 0$, conclui-se que:

$$I(t) = \frac{1}{2}\arctan(t)\ln(1+t^2)$$

Queremos saber o valor de $I(1)$:

$$I(1) = \int_0^1\frac{\ln(x+1)}{x^2+1}dx = \frac{\pi}{8}\ln(2)$$

Como queríamos demonstrar.
\newline \newline
\textbf{Solução meme (by Libardi):}
\newline
Defina a função $I:\mathbb{R}^2\rightarrow\mathbb{R}^2$ de forma que:
$$I(u,v) = \int_0^u\frac{\ln(vx+1)}{x^2+1}dx$$
Trivialmente, temos as condições de contorno $I(0,v) = 0, \forall v\in\mathbb{R}$ e $I(u,0) = 0, \forall u\in\mathbb{R}$. Ademais, temos que as derivadas parciais de primeira ordem de $I$ são:

$$\frac{\partial I}{\partial u} = \frac{\ln(uv+1)}{u^2+1}$$

$$\frac{\partial I}{\partial v} = \frac{\ln(u^2+1)+2v\arctan(u)-2\ln(uv+1)}{2(1+v^2)}$$

Com isso, montamos a equação diferencial parcial:

$$(1+u^2)\frac{\partial I}{\partial u}+(1+v^2)\frac{\partial I}{\partial v} = \frac{1}{2}\ln(u^2+1)+v\arctan(v)$$

Para resolvê-la, performa-se a transformada $(u,v)\rightarrow(\xi,\omega)$ tal que $\xi = u$ e $\omega = \arctan(v)-\arctan(u)$, com a qual se obtém:

$$(1+\xi^2)\frac{\partial I}{\partial \xi} = \frac{1}{2}\ln(\xi^2+1)+\tan(\omega+\arctan(\xi))\arctan(\xi)$$

$$\Rightarrow I = \int_0^u\left(\frac{\ln(\xi^2+1)+2\tan(\omega+\arctan(\xi))\arctan(\xi)}{2(\xi^2+1)}\right) d\xi$$

Haja vista que $\int\frac{\ln(\xi^2+1)}{\xi^2+1}d\xi = \arctan(\xi)\ln(\xi^2+1)-2\int\frac{\xi\arctan(\xi)}{\xi^2+1}d\xi$, temos:

$$I = \frac{1}{2}\arctan(u)\ln(u^2+1)+\int_0^u\frac{\arctan(\xi)}{\xi^2+1}\left(\frac{\tan(\omega)+\xi}{1-\tan(\omega)\xi} - \xi\right)d\xi$$

$$\Leftrightarrow\hspace{0.5cm} I=\frac{1}{2}\arctan(u)\ln(u^2+1)+\tan(\omega)\int_0^u\frac{\arctan(\xi)}{1-\tan(\omega)\xi}d\xi$$

Finalmente, sabendo que $\tan(\omega)\int\frac{\arctan(\xi)}{1-\tan(\omega)\xi}d\xi=-\arctan(\xi)\ln(1-\tan(\omega)\xi)+\int\frac{(\ln(1-\tan(\omega)\xi)}{\xi^2+1}d\xi$, concluímos o problema:

$$\Rightarrow I=\frac{1}{2}\arctan(u)\ln(u^2+1)-\arctan(u)\ln(1-\tan(\omega)u) + \int_0^u\frac{\ln(1-\tan(\omega)\xi)}{1+\xi^2}d\xi$$

$$\Leftrightarrow\; I(u,v)=\frac{1}{2}\arctan(u)\ln(u^2+1)-\arctan(u)\ln\left(1-\frac{u-v}{1+uv}u\right) + I\left(u,\frac{u-v}{1+uv}\right)$$

pois $\tan(\omega)=\frac{v-u}{1+uv}$, e agora simplifica-se a expressão:

$$\Leftrightarrow\; I(u,v)=\arctan(u)\ln\left(\frac{1+uv}{\sqrt{1+u^2}}\right)+I\left(u,\frac{u-v}{1+uv}\right)$$

Substituindo $u=v=1$, calcula-se a integral desejada:

$$I(1,1) = \int_0^1\frac{\ln(x+1)}{x^2+1}dx = \frac{\pi}{8}\ln(2)$$

Como queríamos demonstrar.



\section{Revista do Professor de Matemática 103 problema 454}

\begin{center}
\begin{tikzpicture}[scale=3]

\draw (-1,0) node[left] {$C$}-- (1,0) node[right] {$B$} -- (0, 1.73205) node[above] {$A$}-- (-1,0);
\draw (0,1.09637) circle (0.317837);
\draw (0,0.31783) circle (0.317837);
\draw (-1,0) -- (0.42020,1.0042);
\draw (1,0) -- (-0.42020,1.0042);

\draw (-1,0) -- (0,1.09637) node[above right] {$O_1$} -- (0, 1.73205);
\draw (-1,0) -- (0,0.31783) node[above right] {$O_2$} -- (0,0);

\end{tikzpicture}
\end{center}

Dando os nomes $\alpha$ ao angulo $O_2CB$ e $\beta$ para $ACO_1$, temos que $\alpha+\beta=30^\circ$. Chame de $P_1$ e $P_2$ os pés das alturas dos triângulos $ACO_1$ e $O_2CB$ relativas aos lados $AC$ e $CB$ e $r$ o raio das esferas.

Pelo triângulo $CO_2P_2$:

$$\tan(\alpha)=r$$

e como o ângulo $AP_1O_1$ é $30^\circ$, a distância $AP_1$ é $r\sqrt{3}$, então no triangulo $CP_1O_1$

$$\tan(\beta)=\frac{r}{2-r\sqrt{3}}.$$

Agora basta resolver as equações. Temos que

$$\tan\left(\frac\pi 6-x\right)=\frac{1-\sqrt{3}\tan(x)}{\sqrt{3}+\tan(x)}$$

então

$$\frac{1-r\sqrt{3}}{\sqrt{3}+r}=\frac{r}{2-r\sqrt{3}}$$

$$r=\sqrt{3}-\sqrt{2}.$$

\section{Putnam and Beyond 1.1.3 (generalizado)}

\section{Revista do Professor de Matemática 103 problema 453/Dissertação lokona (generalizado)}

Nomeando as raízes do polinômio $x_1,x_2\dots x_{2021}$, por uma das relações de Viète temos

$$\sum_{h=1}^{2021}x_h=12126$$

$$\frac 1{2021}\sum_{h=1}^{2021}x_h=6$$

Agora trocando as variáveis for $y_i=7-x_i$. Suponha agora que o polinômio não tem raízes maiores que 7, então todos os $y_i$ são positivos.

$$\frac 1{2021}\sum_{h=1}^{2021}y_h=1$$

e

$$P(7)=\prod_{h=1}^{2021}y_h$$

$$P(7)^{1/2021}=\prod_{h=1}^{2021}y_h^{1/2021}$$

Então pela desigualdade das médias aritmética e geométrica

$$1\geq P(7)^{1/2021}$$

$$P(7)\leq 1$$

Contradizendo o enunciado. Portanto nossa suposição que todas as raízes são menores que 7 é falsa.

\section{Revista do Professor de Matemática 103 problema 452 (generalizado)}

Definindo $n+1=a^2$, $\frac n2 +1=b^2$ temos

$$a^2-2b^2=-1$$

que é uma equação de Pell negativa. Suas soluções são dadas por

$$a+b\sqrt{2}=(1+\sqrt{2})^k$$

onde $k$ é um inteiro impar. 

\part{Apêndice}

\chapter{Sugestões de Materiais}

\section{Meta}

O template usado é \href{https://www.latextemplates.com/template/the-legrand-orange-book}{este aqui}. Para saber mais sobre \LaTeX, os \href{https://www.overleaf.com/learn}{materiais do overleaf} tem muita informação e \href{https://tex.stackexchange.com/}{o stack exchange de tex} é bom para tirar duvidas. 

\section{Sites de discussão}

O \href{https://math.stackexchange.com/}{stack exchange de matemática} e sua versão mais avançada, \href{https://mathoverflow.net/}{o mathoverflow} são ótimos para fazer e responder perguntas.

\section{Funções geratrizes}\index{Funções geratrizes}

Esse livro\cite{gfology} tem cinco vezes tudo o que você jamais quis saber sobre funções geratrizes. Ele está disponível grátis em PDF no site do autor.

\section{Bibliografia OBM}

O site da OBM tem uma \href{https://www.obm.org.br/como-se-preparar/bibliografia/}{página com links} sobre matemática e competições de matemática.

\section{SiM}

O ICMC organiza um simpósio anual sobre matemática chamado SiM. \href{https://sites.google.com/usp.br/sim-2020/pagina-inicial?authuser=1}{Aqui} está o link para o de 2020, e para \href{https://sites.google.com/usp.br/sim-2020/eventos/problemas-matematicos?authuser=1}{a lista de problemas} que foram sugeridos esse ano.

%----------------------------------------------------------------------------------------
%	BIBLIOGRAPHY
%----------------------------------------------------------------------------------------

\chapter*{Bibliografia}
\addcontentsline{toc}{chapter}{\textcolor{ocre}{Bibliografia}} % Add a Bibliography heading to the table of contents

%------------------------------------------------

\section*{Artigos}
\addcontentsline{toc}{section}{Articles}
\printbibliography[heading=bibempty,type=article]

%------------------------------------------------

\section*{Livros}
\addcontentsline{toc}{section}{Books}
\printbibliography[heading=bibempty,type=book]

%----------------------------------------------------------------------------------------
%	INDEX
%----------------------------------------------------------------------------------------

\cleardoublepage % Make sure the index starts on an odd (right side) page
\phantomsection
\setlength{\columnsep}{0.75cm} % Space between the 2 columns of the index
\addcontentsline{toc}{chapter}{\textcolor{ocre}{Índice Remissivo}} % Add an Index heading to the table of contents
\printindex % Output the index

%----------------------------------------------------------------------------------------

\end{document}
